% ============================================================================
% Chapter 6 — Conclusion and Future Work (rewritten 2026-10-19 on the final
% numbers: baseline re-run under the comparison protocol, one scoring rule,
% model and delivered views). Numbers trace to the chapter-five tables.
% ============================================================================

\chapter{Conclusion and Future Work}
\label{ch:conclusion}

This chapter states what the thesis built and what its evaluation showed, judges the four hypotheses, collects the limitations, and turns them into future work.

\section{Conclusion}
\label{sec:conclusion}

\subsection{Contributions}
\label{sec:concl-contributions}

The thesis set out to reduce confident hallucination in factual question answering on a small open generator by changing the system around the model rather than the model alone (\Cref{ch:introduction}). It contributes the following.
\begin{itemize}
    \item \textbf{An architecture:} CAEM routes each question to a verified episodic memory, to zero-shot generation or to retrieval-augmented generation, with a retrieval-utility classifier that decides when retrieval is likely to help (\Cref{ch:methodology}). A verifier of nine signals produces a calibrated probability of correctness; a four-outcome decision stores, defers, abstains on or discards each answer; stored answers are re-checked at every cycle boundary and used to fine-tune a low-rank adapter under a retention guard that rolls the adapter back while keeping the memory.
    \item \textbf{A per-answer reliability signal:} every answer reaches the user with its probability and a status tag (verified, provisional, conflicting or insufficient evidence), and the answer is withheld only when the evidence is insufficient (\Cref{sec:decision}).
    \item \textbf{A formal analysis of the stored pool:} two short proven results, a testable precondition and a design target: that storage makes the pool purer than the generator's raw output exactly when the verifier's balanced accuracy exceeds one half (\Cref{thm:bayes-purity}) and that the pool's precision is bounded below by the calibration precision when new questions resemble the calibration data (\Cref{thm:purity}), together with design heuristics whose idealising assumptions are stated and checked against the run (\Cref{sec:theory}).
    \item \textbf{A matched evaluation:} a six-cycle run on five benchmarks against eight baselines under one comparison protocol (the same questions, prompt, scoring rule and hallucination instrument for every system), reported in a model view and a delivered view, with paired tests corrected for multiple comparisons and every reported number rebuilt from the saved answers by one command (\Cref{sec:setup}, \Cref{app:reproducibility}).
    \item \textbf{A diagnosed limit:} a measured account of where verifier-gated self-improvement helps, where it stops helping and why (\Cref{sec:discussion}).
\end{itemize}

\subsection{Headline finding}
\label{sec:concl-headline}

CAEM is a label-efficient system that matches strong baselines, gives users a status tag that genuinely predicts correctness, and shows a measured, diagnosed limit of verifier-gated self-training. The evidence for each part of that statement is as follows.

\textbf{Accuracy and hallucination:} the architecture alone added almost no accuracy: pooled exact match was $0.517$ at C0 against $0.511$ for the zero-shot baseline B1, and CHM was higher ($0.155$ against $0.132$) because the untrained system refused too often (\Cref{tab:headline-steps}). One or two cycles of self-improvement removed most of the over-refusal and gave the best readings of the run ($0.535$ and $0.127$ at C2); later cycles eroded both, and the run ended level with B1 ($0.508$ and $0.134$ at C5). No difference between CAEM and B1 is significant on any benchmark at C2 or C5.

\textbf{Against the baselines:} under the matched protocol CAEM is level with the strong prompting baselines and semantic-entropy abstention (B1, B2, B5, B8), with no significant exact-match win against any of them; it is significantly better than the retrieval baselines and FLARE on several benchmarks; and it matches the exact match of vanilla fine-tuning (B6, $0.512$), which used a gold answer for each of $23{,}500$ training questions against about $2{,}100$ gold labels for CAEM, while B6 reached the lower CHM ($0.103$) (\Cref{tab:baseline-pooled,tab:sig-summary}). An earlier comparison that gave the baselines a shorter prompt had shown CAEM ahead on both measures; the matched comparison removed most of that lead.

\textbf{Selectivity:} what the architecture measurably adds is a reliable signal of which answers to trust. Answers tagged verified were right $65$ to $75$ percent of the time at every cycle, against $51$ percent for B1 over all questions, and the tags are ordered by accuracy, verified above provisional above the other two (\Cref{tab:status-tags}). At C0 CAEM was wrong on $40.5$ percent of the questions it answered, less often than any baseline.

\textbf{Where it works:} the gains follow the verifier's ability to check answers against retrieved passages. On TriviaQA, where that ability is highest (area under the ROC curve $0.75$ to $0.86$), CAEM is most accurate relative to B1; on CommonsenseQA, StrategyQA and TruthfulQA, where the verifier is close to chance, it is level with or below the strong baselines (\Cref{sec:disc-verifier-scope}).

\textbf{Why self-improvement stopped helping:} the stored pool stayed about three-quarters correct, as the theory predicts, but the generator trained on it. Training on answers a verifier has passed converges toward the verifier's own view, as Yi et al.~\cite{fu2025collapse} show: confident confabulation rose from $0.208$ for B1 to $0.305$ at C5, so the verifier scored more of the fine-tuned generator's wrong answers highly. Training on chains written under the retrieval prompt also taught the generator to cite a context it was not given, in $93$ to $99$ percent of zero-shot answers after the first fine-tune (\Cref{sec:disc-phantom}).

\textbf{Routing:} the retrieval-utility classifier, with accompanying corrections, raised pooled exact match by $4.9$ points at cycle three, mainly by retrieving less on questions where retrieval misleads (\Cref{sec:ruc-ablation}).

\subsection{Hypothesis verdicts}
\label{sec:concl-hypotheses}

The four pre-specified hypotheses of \Cref{sec:hypotheses} are judged in \Cref{tab:hypothesis-verdicts}.
\begin{itemize}
    \item \textbf{H1, supported:} the gate's balanced accuracy on the calibration fold stayed above one half at every cycle ($0.615$ to $0.664$ pooled), so the precondition for purification held; CommonsenseQA was the weakest benchmark ($0.52$ to $0.55$).
    \item \textbf{H2, unsupported:} pool purity did not rise from cycle to cycle; calibration-fold precision fell at C1, C2 and C5, within sampling noise.
    \item \textbf{H3, supported:} precision stayed in a band of width $0.014$ ($0.732$ to $0.746$) over the kept cycles; a convergence rate cannot be estimated from six cycles.
    \item \textbf{H4, partially supported:} CHM fell after the first cycle and stayed between $0.127$ and $0.137$, consistent with flattening, but the limit itself cannot be estimated.
\end{itemize}
The two proven results concern the stored pool, and both held where their assumptions held. None of the formal statements bounds the accuracy of the model trained on the pool, which is where the run fell short.

\subsection{Limitations}
\label{sec:concl-open}

The limitations are discussed with their evidence in \Cref{sec:disc-threats,sec:disc-limitations}; the ones that bound the conclusions above are these.
\begin{itemize}
    \item \textbf{One seed, one backbone, five benchmarks:} every result comes from a single run of a three-billion-parameter generator on $300$ evaluation questions per benchmark, and only the paired tests account for sampling noise.
    \item \textbf{Six cycles:} the run stopped after cycle five, as planned after the cycle-four rollback, which is too short to estimate convergence rates or long-run limits.
    \item \textbf{Verifier scope:} the verifier separates correct from incorrect answers well only where passages can confirm an answer, and its selectivity weakened as the generator it judged changed.
    \item \textbf{The hallucination instrument:} two of the seven CHM subtypes use CAEM's own verifier; exact match and the delivered view's error among answered do not depend on it.
    \item \textbf{No repeated questions:} the evaluation used fresh questions to measure learning, so the memory-direct tier and its cost saving were tested only by the router probe (\Cref{tab:tier1-probe}).
    \item \textbf{Bundled ablation:} the classifier ablation also carries pipeline corrections, so it measures the classifier together with them.
    \item \textbf{Labels:} CAEM is label-efficient, not label-free; it used about $2{,}100$ gold labels, and a deployment needs a labelled batch at every cycle (\Cref{sec:conformal-cycle-recalib}).
\end{itemize}

\subsection{Closing remark}
\label{sec:concl-closing}

The thesis began from the observation that a user cannot tell a confident wrong answer from a confident right one. CAEM does not make a small generator much more accurate, and its self-improvement loop helped for only one or two cycles. What it does provide is an answer that carries an honest estimate of its own reliability, a pool of stored answers that stayed as pure as the theory's floor predicts, and a measured account of why training on verifier-passed answers stopped helping, consistent with the verifier's blind spots becoming the generator's. The most useful next step is therefore a better verifier rather than more cycles.

\section{Future Work}
\label{sec:future-work}

The items below follow from the limitations above and from the points that \Cref{ch:methodology,ch:evaluation} left open. The first three address the causes of the limit the run measured; the rest complete measurements this study did not take.

\subsection{A verifier for answers that passages cannot settle}
\label{sec:future-judge}

The verifier is close to chance on commonsense, multi-step and misconception questions because its informative signals compare the answer with passages or with the question (\Cref{sec:disc-verifier-scope}). Four changes follow.
\begin{itemize}
    \item \textbf{Answer-level internal signals:} the token signal scores the whole response, so it measures fluency more than confidence in the answer (\Cref{sec:verifier-internal}); scoring the answer line alone, and computing semantic entropy over the final answers of sampled responses, the setting in which Farquhar et al.~\cite{farquhar2024semantic} found it effective, would give internal signals that track correctness.
    \item \textbf{Checking the reasoning:} every signal scores the final answer, so a chain with a wrong step and a correct conclusion passes. The atomic-fact check gave the same value as the passage-mean signal on every answer of the run, because the answers were too short to decompose (\Cref{sec:verifier-grounding}); applying it to the reasoning chain, where several facts are asserted, would give it something to check and would catch such chains.
    \item \textbf{A learned verifier:} the composite combines fixed signals; a verifier trained on the calibration fold to judge answer and reasoning together, as V-STaR~\cite{hosseini2024vstar} trains its verifier on correct and incorrect self-generated solutions, could learn cues the fixed signals miss. The long-answer judge that was dropped at cycle zero (\Cref{sec:adj-judge-ablation}) could be revisited in the same way.
    \item \textbf{Off-topic answers:} the question-match signal acts as a correctness cue on short answers and does not detect off-topic ones (\Cref{sec:verifier-qa-rel}); a check on the full response, or a judge asked directly whether the answer addresses the question, would let the off-topic subtype of CHM be measured.
\end{itemize}

\subsection{Contradiction-aware verification}
\label{sec:future-refutation}

The entailment judge returns only whether a passage supports a claim, so an answer that a passage contradicts scores the same as one the passages ignore, and the factual-contradiction subtype of CHM could not be measured (\Cref{sec:setup-metrics}). A three-way judge (support, contradiction, neither) would let the decision discard contradicted answers outright and would complete the metric. The language-model judge tried for the two unmeasured subtypes agreed too weakly with hand labels (Cohen's kappa about $0.35$), so a better-validated judge is needed first.

\subsection{Cleaner self-improvement}
\label{sec:future-lora}

The loop's gains faded after two cycles while confident confabulation rose. Four changes target the cause.
\begin{itemize}
    \item \textbf{Remove context references from training chains:} rewriting or removing references to a context in retrieval-written chains before fine-tuning would have prevented the phantom-context citations (\Cref{sec:disc-phantom}).
    \item \textbf{Weight training answers by confidence:} about one stored answer in four is wrong, and the wrong ones carry lower scores; weighting each training answer by its calibrated probability, or raising the training threshold as the pool grows, would reduce the noise the generator learns from.
    \item \textbf{Test the cause directly:} retraining from the same starting point on a pool with the wrong answers removed would show whether the erosion after C2 comes from the wrong answers in the pool, as the account above suggests, or from repeated fine-tuning itself.
    \item \textbf{Retry when unsure:} answers the verifier is unsure about are now deferred or withheld; regenerating them, with retrieval or with more samples, before deciding would turn some of those into verified answers.
    \item \textbf{Adapter capacity:} Biderman et al.~\cite{biderman2024lora} find that low-rank adaptation forgets less but learns less than full fine-tuning; comparing ranks, or full fine-tuning on the verified pool, would show whether the plateau is set by the adapter or by the pool.
\end{itemize}

\subsection{Scope to checkable questions}
\label{sec:future-transfer}

The other direction the verifier result points to is to narrow the system rather than widen the verifier: to deploy CAEM only where answers can be checked against documents, such as trivia, entity and document-grounded questions, and to measure it on a panel of such benchmarks. Because the main run fine-tuned on all three training benchmarks, such a study needs its own run rather than a re-reading of this one. Whether a narrow system or a general verifier is the better path is left open.

\subsection{Isolating the effects}
\label{sec:future-ablations}

\begin{itemize}
    \item \textbf{The classifier alone:} a run with the pipeline corrections but without the retrieval-utility classifier would separate the two parts of the measured $4.9$-point gain (\Cref{sec:ruc-ablation}).
    \item \textbf{Multiple seeds:} repeating the run with different seeds would show how much of the cycle-to-cycle movement, of a size similar to the effects measured, is noise.
    \item \textbf{A flexible-match re-scoring:} part of the early gain may be format learning (\Cref{sec:disc-limitations}); scoring with a flexible answer match as well as exact match would separate the two.
\end{itemize}

\subsection{Testing the formal heuristics}
\label{sec:future-fixed-composite}

The monotonicity heuristic and hypothesis H2 assume a fixed verifier, while the run refitted the composite every cycle. A cycle in which the composite is frozen and only the generator changes would separate the two effects (\Cref{sec:theory}). Measuring the gate's true- and false-positive rates and the generator's accuracy on a labelled sample at every cycle would check the purification condition directly rather than through balanced accuracy.

\subsection{A longer horizon}
\label{sec:future-equilibrium-horizon}

A convergence rate, the long-run limit of CHM, and the equilibrium behaviour under rollbacks all need at least fifteen kept cycles under a fixed composite (\Cref{sec:theory}). A longer run would also show whether the deferred buffer, which ran only three reconsideration passes, recovers more answers over time, and whether a rotating calibration fold behaves better than a fixed one over many cycles (\Cref{sec:conformal-cycle-recalib}).

\subsection{Measurements this study did not take}
\label{sec:future-latency}

\begin{itemize}
    \item \textbf{Time per tier:} the batched evaluation log could not separate the time spent in each tier; timing each tier at a batch size of one would fill in the deployment-cost model (\Cref{sec:tier-cost}).
    \item \textbf{The corpus ceiling:} checking, for every evaluation question, whether any retrieved passage contains the gold answer would bound what retrieval can achieve (\Cref{cor:corpus-floor}).
    \item \textbf{Calibration on more data:} the per-benchmark calibration curves are fitted on $500$ answers, below the size at which Niculescu-Mizil and Caruana~\cite{niculescu2005predicting} find isotonic regression reliable; a larger calibration fold, or a comparison with a parametric fit per benchmark, would show how much the shrinkage toward the pooled fit leaves to chance.
    \item \textbf{Retention of the fine-tuned baseline:} the retention probes were run only on CAEM; running them on B6 would show whether the retention guard protects anything that plain fine-tuning loses.
\end{itemize}

\subsection{Larger backbones and other domains}
\label{sec:future-backbone}

All results come from one three-billion-parameter generator on English factual question answering. Song et al.~\cite{song2024gap} find that small models often gain little from verifying their own outputs, so repeating the study on a larger generator, and on other domains and languages, would show how far the findings carry.

\subsection{Deployment study}
\label{sec:future-deployment}

The memory-direct tier and the production design (\Cref{sec:conformal-cycle-recalib}) were built for traffic in which users repeat questions, which this study did not have. A deployment with live users would measure the share of repeated questions, the cost saving of the memory-direct tier, how users read the status tag, and the labelling effort a cycle needs in practice. It would also test the memory at a realistic size: the run's memory reached $12{,}175$ entries, so the approximate search and the value-based pruning built for millions of entries were never used.
